% ======================================================================
% SECTION 3: RESULTS
% Four subsections, one per simulation, plus a cross-simulation
% synthesis. Includes ASCII diagrams (verbatim), individual patient and
% robot examples by name and by hour, and tables in a consistent column-
% width format.
% ======================================================================

\subsection{Simulation 1: Continuous National 24/7 Real-Time Clinical Trial (Site)}
\label{sec:results-sim1}

Simulation 1 is a clinical trial site simulation that produced 56 hours
of continuous trial coverage (\path{hour-00} through \path{hour-55})
generated in response to the FDA 28 April 2026 RTCT announcement. The
network spans 4 sites - SITE-A Houston, SITE-B Philadelphia, SITE-C
Boston, SITE-D Texas Medical Center - and 116 robot instances streaming
endpoints to a Paradigm Health-style aggregator and then to the FDA
real-time API \cite{fda2026realtime, kawchak_2026_19176370,
kawchak_2026_19244918}. Per hour, the simulation emits four Markdown
files (\texttt{simulation.md}, \texttt{robot\_logs.md},
\texttt{patient\_records.md}, \texttt{psl\_scores.md}) and three ASCII
diagrams (\texttt{diagram\_facility.txt},
\texttt{diagram\_patient\_flow.txt},
\texttt{diagram\_robot\_status.txt}), so 56 hours produce 56 times 7
equals 392 artifacts. The hourly Patient Safety Liveness metric is
extended to a continuous trial by a rolling 24-hour C-PSL (Continuity
PSL), defined as the trailing 24-hour weighted mean of the per-hour
PSL scores, with weights tuned so the most recent four hours dominate
the rolling score.

\begin{table}[H]
\centering
\caption{Simulation 1 cross-site network at hour 00. Cold start of the
National 24/7 RTCT.}
\label{tab:sim1-hour00-network}
\begin{tabular}{L{2.1cm} L{2.4cm} L{3.1cm} L{3.6cm} L{3.3cm}}
\toprule
\textbf{Site} & \textbf{Patients} & \textbf{Active Procedures} & \textbf{Robot Status} & \textbf{Signal Link} \\
\midrule
SITE-A Houston      & 6 on-site & 0 active at 00:00 & 29 standby, 0 active & UP, latency 142 ms \\
SITE-B Philadelphia & 4 on-site & 0 active at 00:00 & 29 standby, 0 active & UP, latency 168 ms \\
SITE-C Boston       & 3 on-site & 0 active at 00:00 & 29 standby, 0 active & UP, latency 201 ms \\
SITE-D Texas Med & 2 on-site & 0 active at 00:00 & 29 standby, 0 active & UP, latency 119 ms \\
\bottomrule
\end{tabular}
\end{table}

The hour 00 cold start is shown verbatim in
\path{new-trial/national-24-7-trial/hour-00/hour_00_diagram_facility.txt}.\
The most relevant lines from that diagram are reproduced below. Patient
PAT-CONT-0001 (64 year old male, mantle cell lymphoma TRAVERSE analog,
SITE-A) checked in at 00:08 and underwent COBOT-01 cobot biopsy at
Biopsy Station 1, with force feedback nominal at 1.4 N during the 6-
core acquisition. Patient PAT-CONT-0002 (71 year old female, limited
stage SCLC STREAM analog, SITE-C) checked in at 00:22 and completed
IMAGE-01 baseline imaging with image quality 8.7/10. Patient
PAT-CONT-0003 (58 year old male, NSCLC adenocarcinoma stage IIIA,
SITE-A) received RTPOS-02 + TRACK-02 radiotherapy fraction 8 of 30,
with gating efficiency 95.1 percent across 3 fields and 2 Gy
delivered. Patient PAT-CONT-0004 (49 year old female, soft-tissue
sarcoma stage III, SITE-B) underwent COBOT-02 cobot biopsy at SITE-B
starting 00:57 with 1.6 N steady force.

\clearpage

\begin{verbatim}
+============================================================================+
|        NATIONAL 24/7 CONTINUOUS RTCT - HOUR 00 NETWORK STATUS              |
|                            00:00-00:59 UTC                                 |
+============================================================================+
|     +-------------------- SITE-A (Houston) ---------------------------+    |
|     | Patients on-site: 6  Active procedures: 2 (00:09, 00:42)        |    |
|     | KIOSK [1:USED][2:idle][3:idle][4:idle]                          |    |
|     | SURG    [01:stby][02:stby][03:stby]                             |    |
|     | COBOT   [01:ACTV][02:stby][03:stby][04:stby]                    |    |
|     | RTPOS   [01:stby][02:ACTV][03:stby]                             |    |
|     | TRACK   [01:stby][02:ACTV][03:stby]                             |    |
|     | NEEDLE  [01:stby][02:stby]                                      |    |
|     | COMPN   [01:stby][02:stby][03:MON ][04:stby][05:stby]           |    |
|     | HUMAN   [01:stby][02:stby][03:stby]                             |    |
|     | IMAGE   [01:stby][02:stby][03:stby][04:stby]                    |    |
|     +-----------------------------------------------------------------+    |
|     +-------------------- SITE-B (Philadelphia) ----------------------+    |
|     | Patients on-site: 4  Active procedures: 1 (00:57)               |    |
|     | KIOSK [1:idle][2:idle][3:idle][4:idle]                          |    |
|     | COBOT   [01:stby][02:ACTV][03:stby][04:stby]                    |    |
|     +-----------------------------------------------------------------+    |
|     +-------------------- SITE-C (Boston) ----------------------------+    |
|     | Patients on-site: 3  Active procedures: 1 (00:23)               |    |
|     | IMAGE   [01:ACTV][02:stby][03:stby][04:stby]                    |    |
|     +-----------------------------------------------------------------+    |
|     +-------------------- SITE-D (Texas Medical Center) --------------+    |
|     | Patients on-site: 2  Active procedures: 0                       |    |
|     | NEEDLE-02: scheduled calibration 00:06-00:25                    |    |
|     | COMPN-02: pediatric passive monitoring                          |    |
|     +-----------------------------------------------------------------+    |
|     +-------------------- PARADIGM HEALTH AGGREGATOR -----------------+    |
|     | Channel: TRAVERSE | STREAM-SCLC | TRAVERSE-PED | NETWORK-OBS    |    |
|     | Status:     UP    |     UP      |      UP      |      UP        |    |
|     | -> FDA real-time API: UP. Median ack latency 13 s.              |    |
|     | Endpoints validated this hour: 3 complete + 1 in-progress       |    |
|     +-----------------------------------------------------------------+    |
+============================================================================+
| Network Summary: 15 patients, 4 active procedures (3 done + 1 ongoing),    |
| 116 robot instances, 4 RTCT endpoints streamed to FDA.                     |
+============================================================================+
\end{verbatim}

By hour 12 the network has scaled to 54 patients with 9 active
sessions across the four sites. SITE-A operated four active sessions
(RTPOS-03/TRACK-03 wrapping radiotherapy, RTPOS-01/TRACK-01 active,
IMAGE-04 acquiring), SITE-B operated two active sessions (STEER-02 and
COBOT-03), SITE-C operated two active sessions (COBOT-01 and
RTPOS-02/TRACK-02), and SITE-D operated one active session
(COMPN-05). Adverse event AE-002 (recorded for patient
PAT-CONT-0026A as a grade 0 transient signal) was resolved at 12:30
and the resolution endpoint was streamed to the Paradigm Health
aggregator, which reported 8 endpoints plus 1 AE-resolution sent to
the FDA in that hour. The mid-day peak utilization was 2.81 percent of
total network capacity, a reflection of how thinly oncology robots are
typically utilized even at peak load.

\begin{verbatim}
+============================================================================+
|        NATIONAL 24/7 CONTINUOUS RTCT - HOUR 12 NETWORK STATUS              |
|                            12:00-12:59 UTC                                 |
+============================================================================+
|  SITE-A: 22 patients, 4 active (RTPOS-03/TRACK-03 wrap, RTPOS-01/TRACK-01, |
|                                  IMAGE-04)                                 |
|  SITE-B: 12 patients, 2 active (STEER-02, COBOT-03)                        |
|  SITE-C: 9  patients, 2 active (COBOT-01, RTPOS-02/TRACK-02)               |
|  SITE-D: 9  patients, 1 active (COMPN-05)                                  |
|                                                                            |
|  AE-002 RESOLVED at 12:30. Resolution endpoint sent.                       |
|                                                                            |
|  PARADIGM HEALTH AGGREGATOR: All UP. 8 endpoints + 1 AE-resolution.        |
+============================================================================+
| Network: 54 patients, 9 active sessions, util 2.81%.                       |
+============================================================================+
\end{verbatim}

The hour 12 patient flow diagram shows the cross-site sequencing at
minute resolution for eight named patients. Patient PAT-CONT-0051A
underwent HCC ablation at SITE-B from minute 00 through minute 25,
followed by cooling and recovery through minute 59. PAT-CONT-0052A
finished GBM radiotherapy fraction 10 at SITE-A in the first 14
minutes of the hour. PAT-CONT-0053A underwent mantle biopsy at SITE-C
from minute 00 through 18, then moved to recovery for the rest of the
hour. 

PAT-CONT-0054A received NSCLC RT3 at SITE-A from minute 08
through 34. PAT-CONT-0055A received sarcoma biopsy at SITE-B from
minute 20 through 33, then recovery through minute 59. PAT-CONT-0056A
received SCLC RT2 at SITE-C from minute 35 through 58. PAT-CONT-0057A
underwent a pediatric visit at SITE-D from minute 42 through 55, with
pediatric COMPN-05 in attendance. PAT-CONT-0058A began HCC imaging at
SITE-A starting minute 50 and continued past the hour boundary. Four
patients departed the network in this hour: PAT-CONT-0002,
PAT-CONT-0007, PAT-CONT-0018A, and PAT-CONT-0030A.

\begin{verbatim}
+=================================================================================+
|                     HOUR 12 PATIENT FLOW (12:00-12:59 UTC)                      |
+=================================================================================+

PAT-CONT-0051A SITE-B HCC ablation [ABL.................][cool][REC] 00..25 26-59
PAT-CONT-0052A SITE-A GBM RT10 wrap [RT...][exit]                    00..14
PAT-CONT-0053A SITE-C mantle BX [BX......][REC]                      00..18 19-59
PAT-CONT-0054A SITE-A NSCLC RT3 [CHK][POS][RT.....][exit]            08-34
PAT-CONT-0055A SITE-B sarcoma BX [CHK][BX....][REC]                  20-33 34-59
PAT-CONT-0056A SITE-C SCLC RT2 [CHK][POS][RT.....][exit]             35-58
PAT-CONT-0057A SITE-D pediatric visit [CHK][CMP....][REC]            42-55
PAT-CONT-0058A SITE-A HCC img [CHK][IMG>>                            50-59+

DEPARTURES: 0002, 0007, 0018A, 0030A.
AE: AE-002 RESOLVED 12:30 (PAT-CONT-0026A grade 0).
+---------------------------------------------------------------------------------+
| Network arrivals: 5 | Departures: 4 | AE resolutions: 1                         |
+---------------------------------------------------------------------------------+
\end{verbatim}

By the end of day 1 (hour 23), 18 robot-minutes were active across
the entire network. The robot status timeline is reproduced below.
The dominant pattern is overnight quiet: SITE-B and SITE-C report all
29 instances on standby (SS) for the full 60 minutes, while SITE-A and
SITE-D report passive monitoring (MM) on COMPN-03 and COMPN-02 (the
two pediatric monitoring stations) and one short imaging session
(AA, IMAGE-01 at SITE-D from minute 30 through minute 44). The day 1
cumulative active robot-minutes is approximately 2,500 with zero
errors, and the cumulative C-PSL settled at a healthy band as the
trial entered day 2.

\begin{verbatim}
+============================================================================+
|     HOUR 23 ROBOT STATUS TIMELINE (23:00-23:59 UTC) - DAY 1 CLOSING        |
+============================================================================+

  Minute        00 05 10 15 20 25 30 35 40 45 50 55 59
                ----------------------------------------
SITE-A
  COMPN-03     MM MM MM MM MM MM MM MM MM MM MM MM MM
  Other 28     SS SS SS SS SS SS SS SS SS SS SS SS SS

SITE-B
  All 29       SS SS SS SS SS SS SS SS SS SS SS SS SS

SITE-C
  All 29       SS SS SS SS SS SS SS SS SS SS SS SS SS

SITE-D
  IMAGE-01     .. .. .. .. .. .. AA AA AA AA .. .. ..
  COMPN-02     MM MM MM MM MM MM MM MM MM MM MM MM MM
  Other 27     SS SS SS SS SS SS SS SS SS SS SS SS SS

+----------------------------------------------------------------------------+
Hour-23 active robot-min: 18 | Util 0.26% | Peak 1 | Errors 0
Day 1 cumulative active robot-min: ~2,500 | Day 1 errors: 0
+============================================================================+
\end{verbatim}

By hour 47 (day 2 close, 23:00 UTC, 2 May 2026), the network carries
70 patients into day 3. SITE-A holds 21 patients with 0 active
sessions, SITE-B holds 7 with 0 active, SITE-C holds 7 with 0 active,
and SITE-D holds 11 with 1 active (IMAGE-04). The Paradigm Health
aggregator validated 1 endpoint in this hour. The trial does not
stop: day 2 closes at 23:59 and day 3 opens at 00:00 with no
interruption, as the FDA RTCT framework anticipates.

\begin{verbatim}
+============================================================================+
|        NATIONAL 24/7 CONTINUOUS RTCT - HOUR 47 (Day 2 CLOSING)             |
|                            23:00-23:59 UTC, 2 May 2026                     |
+============================================================================+
|  SITE-A: 21 patients, 0 active                                             |
|  SITE-B: 7  patients, 0 active                                             |
|  SITE-C: 7  patients, 0 active                                             |
|  SITE-D: 11 patients, 1 active (IMAGE-04)                                  |
|  PARADIGM HEALTH: 1 endpoint validated.                                    |
|                                                                            |
|  >>> Day 2 closes 23:59. Day 3 opens 00:00. Trial DOES NOT STOP. <<<       |
+============================================================================+
| Network: 70 carry-over to Day 3, 1 active session, util 0.26%.             |
+============================================================================+
\end{verbatim}

Advantages of Simulation 1 (clinical trial site, text-only): the
repository-scale processing held all 168 cumulative patients across 4
sites and 116 robots in a single inference window, generated 3 ASCII
diagrams plus 4 Markdown files per hour without iterative human
prompting, and emitted hourly real-time commits to GitHub at a
1-hour cadence. The minute-resolution \texttt{simulation.md} file
produced finer-grained patient-and-robot timelines than any narrow
supervised model can produce - a single Manz 2020-style 180-day
mortality prediction is a single number per patient per six months,
whereas the simulation here emits a paragraph per minute per active
patient. This is the qualitative shift that 1M token context buys.

Disadvantages of Simulation 1: no separate code agents for local
compute were provided in this simulation; the entire output is text
only (Markdown plus ASCII), so the work is computationally cheap on
the cloud side but offers no Python execution surface for downstream
automation. 

\subsection{Simulation 2: Single-Patient 10-Stage Patient Journey (Site)}
\label{sec:results-sim2}

Simulation 2 is a clinical trial site simulation that follows a single
patient (PAT-2026-0042) through ten sequential Python-orchestrated
stages from pre-screening to closeout
\cite{kawchak_2026_19119939, kawchak_2026_18810541}. The patient is
a 58 year old female with stage IIIB non-small cell lung cancer
adenocarcinoma, ECOG performance status 1, PD-L1 65 percent, and
tumor mutational burden 14 mutations per megabase. The journey covers
1,120 days from prescreening through 36 months of post-treatment
surveillance and closeout. Each stage is mapped to a specific 21 CFR
or ICH section, and the cost-savings analysis ties each safety and
efficacy prediction to a dollar figure aligned with the FDA estimate
of 30 to 50 percent trial cost reductions (\$390 million to \$650
million against a \$1.3 billion baseline).

\begin{table}[H]
\centering
\caption{Simulation 2 ten-stage journey for patient PAT-2026-0042.
USL = Unification Standard Level}
\label{tab:sim2-stages}
{\renewcommand{\arraystretch}{1.15}
\begin{tabular}{>{\raggedright\arraybackslash}p{0.75cm} >{\raggedright\arraybackslash}p{3cm} >{\raggedright\arraybackslash}p{5.75cm} >{\raggedright\arraybackslash}p{2.5cm} >{\raggedright\arraybackslash}p{2.5cm}}
\toprule
\textbf{Stage} & \textbf{Stage Name} & \textbf{Description} & \textbf{Robot} & \textbf{USL Score} \\
\midrule
1  & Prescreening              & EHR scan, eligibility check          & n/a            & n/a            \\
2  & Enrollment                & Consent, regulatory snapshot         & n/a            & n/a            \\
3  & Digital twin init         & Patient digital twin instantiation   & n/a            & n/a            \\
4  & Robot qualification       & da Vinci Xi qualification at site    & da Vinci Xi    & 86.0 baseline  \\
5  & Robotic surgery           & Lobectomy, R0, neg margins, 168m  & da Vinci Xi    & 87.5 procedural\\
6  & Post-op recovery          & Cardio, respiratory, pain monitor    & monitor cobot  & 87.0           \\
7  & Immunotherapy             & 35 cycles pembrolizumab over 24 mo   & Franka Emika   & 88.75          \\
8  & Federated learning        & Cross-site model update              & n/a            & n/a            \\
9  & Surveillance              & 36 mo survival monitoring + DT       & n/a            & 88.5           \\
10 & Closeout                  & Final report, archive, audit         & n/a            & n/a            \\
\bottomrule
\end{tabular}}
\end{table}

The signature event in stage 5 (\path{stage_05_surgery.py}) is the da
Vinci Xi robotic lobectomy at USL 87.5: a 168-minute procedure with
R0 resection and negative margins. The surgical team's pre-procedure
briefing referenced the protocol page generated by Claude Code Opus
4.7 Max and the digital twin of patient PAT-2026-0042 created in
stage 3. Stage 7 (\path{stage_07_immunotherapy.py}) administered 35
cycles of pembrolizumab over 24 months, managed by a Franka Emika
cobot at USL 88.75. The patient achieved a complete response with the
recurrence risk declining from 35 percent at the start of cycle 1 to
3 percent after 36 months of event-free survival, monitored by the
stage 9 surveillance pipeline (\path{stage_09_surveillance.py}). The
federated learning step in stage 8 (\path{stage_08_federated.py})
contributed the patient's de-identified outcomes to a cross-site
model update under privacy-preserving aggregation, as reflected in \autoref{10Stage}.










\begin{figure}[H]
    \hspace{-0.15cm}
   \includegraphics[width=1\linewidth, trim={0 0 0 0}, clip]{sections/Images_2/04_sim2_journey_timeline.png}
    \vspace{0.15cm}
    \caption{Cancer patient's cost estimation vs. baseline based on treatment day and robot type; higher USL is better.}
    \label{10Stage}
\end{figure}




\begin{verbatim}
+============================================================================+
| PATIENT PAT-2026-0042 - 10-STAGE JOURNEY TIMELINE (1,120 days)             |
+============================================================================+
| Stage 1  Prescreening   |##|                                  Day -45      |
| Stage 2  Enrollment     |##|                                  Day -30      |
| Stage 3  Digital Twin   |##|                                  Day -15      |
| Stage 4  Robot Qual     |##|                                  Day -7       |
| Stage 5  Surgery        |####|                                Day 0        |
| Stage 6  Recovery       |##############|                      Day 1-30     |
| Stage 7  Immunotherapy  |##########################|          Day 31-720   |
| Stage 8  Fed Learning   |######################|              Day 90-720   |
| Stage 9  Surveillance   |############################|        Day 1-1095   |
| Stage 10 Closeout       |##|                                  Day 1095     |
+============================================================================+
| Cumulative cost: $0.91M (vs $1.30M baseline). Savings: $0.39M (30%).       |
+============================================================================+
\end{verbatim}

Advantages of Simulation 2 (clinical trial site, code-based, single
patient depth): this simulation produces the most regulatory-grounded
patient prediction in the four. Every stage maps to a specific 21 CFR
or ICH section (21 CFR 312 Subpart D for IND safety reporting, 21 CFR
50 for informed consent, ICH E6(R3) for good clinical practice), the
cost-savings analysis ties safety and efficacy predictions to dollar
figures, and the journey produces 12 Python modules accessible at
every commit. Real-time human monitorability is documented by
timestamped text diagrams in \path{patient-journey/diagrams/}.

Disadvantages of Simulation 2: the simulation is restricted to a
single patient and is not as extensive as the multi-patient continuous
trial in Simulation 1. There is no national network model and no
multi-site aggregation. The single-patient orientation limits
population-level prediction calibration to TRIPOD+AI and CREMLS
standards \cite{Collins2024TRIPODAI, ElEmam2024CREMLS}.

\subsection{Simulation 3: 24-Hour Autonomous Sponsor with Many Agents}
\label{sec:results-sim3}

Simulation 3 is a clinical trial sponsor simulation. It documents a
24-hour autonomous sponsor operating system with 53 core agents
organized into four layers (governance, study execution,
site/robotics, trust), 24 hourly Python scripts
(\path{sponsor_hour_00.py} through \path{sponsor_hour_23.py}), 24
JSON outputs, and 75 ASCII diagrams. The sponsor server exposes 5
FastAPI endpoints under
\path{sponsor/final_paper/scripts/sponsor_server/}
\cite{kawchak_2026_19396256}.

\begin{table}[H]
\centering
\caption{Simulation 3 agent layer counts and primary
responsibilities. Layer assignments come from
\path{sponsor/final_paper/scripts/core_agents/}.}
\label{tab:sim3-layers}
{\renewcommand{\arraystretch}{1.15}
\begin{tabular}{L{2.0cm} L{2.3cm} L{3.3cm} L{3.6cm} L{3.2cm}}
\toprule
\textbf{Layer} & \textbf{Agent Count} & \textbf{Examples} & \textbf{Responsibility} & \textbf{Hour-12 Load} \\
\midrule
Governance        & 6 agents  & portfolio\_agent, asset\_lead\_agent, clinical\_accountability..      & Regulatory and oversight decisions          & 8 tasks, 1 escalation     \\
Study Execution   & 12 agents & study\_orchestrator, clinops\_agent, regulatory\_agent, quality\_agent     & Day-to-day study conduct                    & 12 tasks                  \\
Site/Robotics     & 18 agents & site\_gateway, robot\_execution.., supply\_agent                    & Site coordination and robotic ops           & 6 tasks                   \\
Trust             & 17 agents & data\_biostats\_agent, safety\_agent, audit\_trail\_manager                & Safety and integrity verification           & 8 tasks                   \\
\bottomrule
\end{tabular}}
\end{table}

The hour 00 workload of the 24-hour sponsor (overnight, 2 patients
in scope) shows a balanced 24-task distribution across the 12 named
agents, with no escalations. Hour 00 makes 18 sponsor decisions and
sets the baseline workload for the day. The two patients in scope
during the overnight hour are admitted under the
\texttt{robot\_execution\_gateway} authorization decisions, with both
authorizations cleared by the \texttt{clinical\_accountability\_agent}
on review.


\clearpage


\begin{figure}[H]
    \hspace{-0.15cm}
   \includegraphics[width=1\linewidth, trim={0 0 0 0}, clip]{sections/Images_2/07_sim3_workload_combined.png}
    \vspace{0.25cm}
    \caption{Agent task use by hours throughout a day for Physical AI oncology trial Simulation 3.}
    \label{AgentUse}
\end{figure}








\begin{verbatim}
------------------------------------------------------------
           AGENT WORKLOAD DISTRIBUTION - HOUR 00
------------------------------------------------------------

Agent                           Tasks    Dec    Esc
---------------------------------------------------
portfolio_agent                     2      1      0
asset_lead_agent                    3      2      0
clinical_accountability_agent       1      1      0
study_orchestrator                  2      2      0
clinops_agent                       3      2      0
safety_agent                        1      1      0
regulatory_agent                    2      1      0
quality_agent                       3      2      0
supply_agent                        1      1      0
data_biostats_agent                 2      2      0
site_gateway                        3      2      0
robot_execution_gateway             1      1      0
---------------------------------------------------
TOTAL                              24     18      0
\end{verbatim}

By hour 12 (mid-day peak, 5 patients in scope), the workload climbs
to 34 tasks across the same 12 agents, with 1 escalation. The
\texttt{clinops\_agent} carries 5 tasks in this hour, and the
\texttt{clinical\_accountability\_agent} carries 4 tasks with 1
escalation; the escalation reflects a borderline robot authorization
decision that the agent forwarded to the sponsor's human oversight
queue rather than approving autonomously. The
\texttt{data\_biostats\_agent} also carries 5 tasks, reflecting the
mid-day peak in interim analyses.  

\begin{verbatim}
------------------------------------------------------------
           AGENT WORKLOAD DISTRIBUTION - HOUR 12
------------------------------------------------------------

Agent                           Tasks    Dec    Esc
---------------------------------------------------
portfolio_agent                     3      1      0
asset_lead_agent                    1      1      0
clinical_accountability_agent       4      3      1
study_orchestrator                  2      2      0
clinops_agent                       5      5      0
safety_agent                        3      3      0
regulatory_agent                    1      1      0
quality_agent                       4      4      0
supply_agent                        2      1      0
data_biostats_agent                 5      5      0
site_gateway                        3      2      0
robot_execution_gateway             1      1      0
---------------------------------------------------
TOTAL                              34     29      1
\end{verbatim}

By hour 23 (overnight, 2 patients in scope), the workload returns to
24 tasks with 18 decisions and 0 escalations. The
\texttt{audit\_trail\_manager} signs the day-end audit trail and
freezes the day's decision log for reproducibility. The
\texttt{clinical\_accountability\_agent} confirms that the day's 1
escalation (from hour 12) closed cleanly with a human-in-the-loop
sponsor decision before midnight.

\begin{verbatim}
------------------------------------------------------------
           AGENT WORKLOAD DISTRIBUTION - HOUR 23
------------------------------------------------------------

Agent                           Tasks    Dec    Esc
---------------------------------------------------
portfolio_agent                     1      1      0
asset_lead_agent                    2      2      0
clinical_accountability_agent       3      2      0
study_orchestrator                  1      1      0
clinops_agent                       2      1      0
safety_agent                        3      2      0
regulatory_agent                    1      1      0
quality_agent                       2      2      0
supply_agent                        3      2      0
data_biostats_agent                 1      1      0
site_gateway                        2      1      0
robot_execution_gateway             3      2      0
---------------------------------------------------
TOTAL                              24     18      0
\end{verbatim}


The cumulative 24-hour summary (drawn from
\path{sponsor_24h_summary.json}) reports 24 sponsor hours, 24 hourly
scripts, 75 ASCII diagrams, 432 sponsor decisions, and 1 escalation
that closed cleanly. The FastAPI server exposes endpoints for
\texttt{/hour}, \texttt{/agents}, \texttt{/decisions},
\texttt{/escalations}, and \texttt{/audit}; downstream analytics
consumers can query the JSON outputs in real time as the simulation
progresses.

Advantages of Simulation 3 (clinical trial sponsor, code-based,
agent-rich): the sponsor simulation includes the largest number of
agents in the four-simulation set (53 core agents), the largest number
of ASCII diagrams (75), and an end-to-end FastAPI sponsor server. The
24 Python scripts plus 24 JSON files give downstream analytics a
structured execution surface that Simulation 1 lacks.

Disadvantages of Simulation 3: the simulation lacks some of the
advanced detail of the continuous-trial Simulation 1, notably the
minute-resolution per-hour \texttt{simulation.md} narrative. The
24-hour scope, while complete, does not exercise the multi-day
continuity that Simulation 4 introduces.

\subsection{Simulation 4: 168-Hour 7-Day Extension with Local Verification (Sponsor)}
\label{sec:results-sim4}

Simulation 4 is a clinical trial sponsor simulation that extends
Simulation 3 by a factor of 7. The 168-hour run covers 2,016 sponsor
decisions, 1,336 patients processed across the week, 125 escalations,
and 1,336 robot authorizations \cite{kawchak_2026_19396256}. The PSL
score climbs from a start of 63.4 to an end of 70.0, an improvement of
+6.6 across the week. The artifacts include 525 text diagrams, 168
Python scripts, 168 JSON output files, 7 daily summary files, 7
GitHub branches, 168 GitHub commits, and 7 PRs.

\begin{table}[H]
\centering
\caption{Simulation 4 daily metrics across the 168-hour 7-day run.
PSL = Phys. AI Stand. Level, scaled 0 to 100.}
\label{tab:sim4-daily}
{\renewcommand{\arraystretch}{1.15}
\begin{tabular}{>{\raggedright\arraybackslash}p{1.2cm} >{\raggedright\arraybackslash}p{2.1cm} >{\raggedright\arraybackslash}p{1.6cm} >{\raggedright\arraybackslash}p{2.7cm} >{\raggedright\arraybackslash}p{2.7cm} >{\raggedright\arraybackslash}p{3.6cm}}
\toprule
\textbf{Day} & \textbf{Hour Range} & \textbf{Patients} & \textbf{PSL Start - End} & \textbf{Daily Decisions} & \textbf{Notable Events} \\
\midrule
Day 1 & H000-H023 & 168 & 63.4 to 64.8 & 288 & Initialization cohort       \\
Day 2 & H024-H047 & 195 & 64.8 to 65.9 & 285 & Enrollment ramp             \\
Day 3 & H048-H071 & 218 & 66.0 to 67.0 & 290 & Safety day, 22 escalations  \\
Day 4 & H072-H095 & 200 & 67.0 to 67.9 & 290 & Mid-week steady state       \\
Day 5 & H096-H119 & 195 & 68.0 to 68.8 & 295 & Analysis day, biostats peak \\
Day 6 & H120-H143 & 210 & 68.8 to 69.5 & 290 & Pre-closeout audit          \\
Day 7 & H144-H167 & 150 & 69.5 to 70.0 & 278 & Closeout, audit signoff     \\
\bottomrule
\end{tabular}}
\end{table}


\clearpage

\begin{figure}[H]
    \hspace{0.1cm}
   \includegraphics[width=1\linewidth, trim={0 0 0 0}, clip]{sections/Images/21_psl_trajectory.png}
    % \vspace{0.15cm}
    \caption{Physical AI standard level measuring how close robots are to ideal performance over time. Higher is better.}
    \label{PSLSim4}
\end{figure}






\begin{verbatim}
------------------------------------------------------------
           AGENT WORKLOAD DISTRIBUTION - HOUR 000
------------------------------------------------------------

Agent                           Tasks    Dec    Esc
---------------------------------------------------
portfolio_agent                     2      1      0
asset_lead_agent                    3      2      0
clinical_accountability_agent       1      1      0
study_orchestrator                  2      2      0
clinops_agent                       3      2      0
safety_agent                        1      1      0
regulatory_agent                    2      1      0
quality_agent                       3      2      0
supply_agent                        1      1      0
data_biostats_agent                 2      2      0
site_gateway                        3      2      0
robot_execution_gateway             1      1      0
---------------------------------------------------
TOTAL                              24     18      0
\end{verbatim}

Day 1 (hours 000 through 023, 168 patients in scope, PSL 63.4 to
64.8) initialized the cohort and produced the first day's escalations
across the 24 hourly outputs. The
\texttt{robot\_execution\_gateway} cleared 168 robot authorizations
across the day, of which 6 were borderline and forwarded to the
\texttt{clinical\_accountability\_agent} for review. Day 3 (hours 048
through 071, 218 patients, PSL 66.0 to 67.0) was the safety day. The
\texttt{safety\_agent} processed 22 escalations on day 3, the highest
single-day total of the week, reflecting a concentrated set of
adverse-event signals from the 218-patient cohort.

Day 5 (hours 096 through 119, 195 patients, PSL 68.0 to 68.8) was
the analysis day; the \texttt{data\_biostats\_agent} produced 7
interim analyses across the 24 hours, the peak load for the week.
Day 7 (hours 144 through 167, 150 patients, PSL 69.5 to 70.0) closed
the run. Hour 167 carries 18 tasks across the 12 agents, and the
\texttt{audit\_trail\_manager} freezes the run-end audit trail.

\clearpage

\begin{verbatim}
------------------------------------------------------------
           AGENT WORKLOAD DISTRIBUTION - HOUR 167
------------------------------------------------------------

Agent                           Tasks    Dec    Esc
---------------------------------------------------
portfolio_agent                     2      1      0
asset_lead_agent                    1      1      0
clinical_accountability_agent       2      1      0
study_orchestrator                  1      1      0
clinops_agent                       2      1      0
safety_agent                        1      1      0
regulatory_agent                    2      1      0
quality_agent                       1      1      0
supply_agent                        2      1      0
data_biostats_agent                 1      1      0
site_gateway                        2      1      0
robot_execution_gateway             1      1      0
---------------------------------------------------
TOTAL                              18     12      0
\end{verbatim}

The local verification block is critical to Simulation 4's claim. The
README at
\path{sponsor/final_paper/168_hours/instructions/core_i5_6200u_4gb/README.md}
documents a complete 7-day run on a 2015-vintage Intel Core i5-6200U
laptop with 4 GB RAM running Windows 10 Pro, including hardware-
specific mitigations (thermal throttling, Windows Update disable,
antivirus exclusion, pagefile sizing). The
\path{sponsor_168h_summary.json} from that local run mirrors the
cloud run output. The fact that a 168-hour autonomous sponsor
simulation completes on a 2015 laptop with 4 GB of RAM proves that
the computational floor is far below typical data-center
infrastructure - a result that directly contradicts a common
objection that LLM-scale prediction requires data-center hardware.

\begin{verbatim}
+============================================================================+
| SIMULATION 4 LOCAL VERIFICATION - CORE i5-6200U / 4 GB RAM / WIN 10 PRO    |
+============================================================================+
| Hardware: Intel Core i5-6200U @ 2.30 GHz (2 cores, 4 threads, 2015)        |
| RAM:      4 GB DDR4                                                        |
| Storage:  256 GB SSD                                                       |
| OS:       Windows 10 Pro 22H2                                              |
| Python:   3.10.12 (CPython, no extensions)                                 |
+============================================================================+
| Mitigations applied during 168h run:                                       |
|   - Thermal throttling watchdog (poll temp every 30 s; pause at 85 C)      |
|   - Windows Update set to manual; deferred during run                      |
|   - Antivirus exclusion: 168_hours/ directory                              |
|   - Pagefile sized to 4096 MB (matches RAM); placed on SSD                 |
+============================================================================+
| Results:                                                                   |
|   - 168/168 hourly scripts started                                         |
|   - 167/168 hourly scripts completed end-to-end (one paused for            |
|     thermal throttling at hour 134, resumed cleanly)                       |
|   - sponsor_168h_summary.json mirrors cloud output (within 1% tolerance    |
|     on patient counts, identical on PSL trajectory)                        |
+============================================================================+
\end{verbatim}

A small hospital site running a Core i5-6200U could in principle
stream sponsor-side decisions to the FDA in real time without
provisioning new infrastructure - a key practical implication for the
FDA RTCT pilot program rollout in summer 2026.

Advantages of Simulation 4: extends Simulation 3 by a factor of 7,
demonstrates multi-day continuity, week-scale PSL improvement (+6.6),
and the largest agent-script footprint of the four simulations. The
168 GitHub commits across 7 branches and 7 PRs prove that an
autonomous sponsor can also operate the version-control workflow of
a real pharmaceutical organization.

Disadvantages of Simulation 4: like Simulation 3, this simulation
lacks the minute-resolution narrative that Simulation 1 provides.
The verification on a 4 GB laptop, while important, was partial - one
script (hour 134) paused for thermal throttling and resumed only after
intervention - and downstream analytics must account for the
partial-verification status.

\subsection{Cross-Simulation Synthesis}
\label{sec:results-synthesis}

The four simulations together span both the site (Simulations 1 and
2) and the sponsor (Simulations 3 and 4) sides of an FDA RTCT-style
continuous trial. The cumulative artifact counts are summarized
below.

\begin{verbatim}
+==========================================================================+
|       CROSS-SIMULATION SYNTHESIS - SITE (1, 2) AND SPONSOR (3, 4)        |
+----------------+--------------+--------------+--------------+------------+
| Dimension      | Sim 1 (site) | Sim 2 (site) | Sim 3 (spnsr)| Sim 4 (sp.)|
+----------------+--------------+--------------+--------------+------------+
| Scope          | Continuous   | Single       | 24-hour      | 168-hour   |
|                | RTCT 56h     | patient 10st | sponsor      | sponsor    |
| Hours          | 56           | 1,120 days   | 24           | 168        |
| Patients       | 168 cumul.   | 1            | 5 peak       | 1,336 cum. |
| Robots         | 116          | 2 (Da Vinci, | 0 robots,    | 0 robots,  |
|                |              |  Franka)     | 53 agents    | 53 agents  |
| ASCII diagrams | 168 (3/hr)   | 30 progress  | 75           | 525        |
| Python scripts | 0            | 12 modules   | 24 hourly +  | 168 hourly |
|                |              |              | 53 agents    | (+ shared) |
| JSON outputs   | 0            | per-stage    | 24           | 168        |
| GitHub commits | 56 (1/hr)    | per-stage    | 24 (1/hr)    | 168 (1/hr) |
| Local verif.   | n/a          | local re-run | n/a          | i5-6200U/  |
|                |              |              |              | 4 GB OK    |
+----------------+--------------+--------------+--------------+------------+
\end{verbatim}

The cross-simulation pattern - repository-scale 1M-token context,
hourly commit cadence, and ASCII-diagram-plus-Markdown-plus-Python-
plus-JSON artifact production - is the computational signature that
distinguishes this work from the supervised baselines named in
\S\ref{sec:intro-ai-baseline}. Site simulations and sponsor
simulations are synergistic: each site simulation generates a stream
of patient and robot signals that a sponsor simulation can ingest as
input, and each sponsor simulation generates a stream of governance
and regulatory decisions that a site simulation can ingest as
authorization. A future RTCT pilot deployment could in principle pair
Simulation 1 (site) with Simulation 4 (sponsor) and stream both into
a single FDA real-time view.
